\begin{recipe}
[% 
    preparationtime = {\unit[15]{min}},
    portion = {\portion{1}},
    source = {Adrian Hieber},
    calory = {1110 kcal | 35g Protein | 138g KH | 47g Fett},
    rating = {1}
]
{Nudeln in Erdnusssauce}
    
    % \graph
    % {% pictures
    %     small=pic/erdnussnudeln_small,
    %     big=pic/erdnussnudeln_big
    % }
    
    \introduction{%
        Nichts besonderes oder kompliziertes, einfach etwas was satt macht und nicht nur Nudeln mit Pesto ist.
    }
    
    \ingredients{%
        200 g & Udon-Nudeln\\
        2 Stangen & Frühlingszwiebeln\\
        1 cm & Ingwer\\
        1--2 & Knoblauchzehen\\
        1 EL & Sojasauce\\
        2 EL & Erdnussbutter\\
        1 TL & Chili in Öl\\
        2 TL & Sesamöl\\
        1{,}5 TL & Reisweinessig\\
        1 EL & Sesamkörner
    }
    
    \preparation{%
        \begin{enumerate}
            \item Sesamkörner in einer Pfanne ohne Öl anrösten und beiseitelegen.
            
            \item Ingwer und Knoblauch fein hacken und zusammen mit dem Chiliöl kurz anbraten.
            
            \item Udon-Nudeln in kochendes Wasser geben und nach Packungsangabe (ca. 90--180 Sekunden) garen.
            
            \item Frühlingszwiebeln in Ringe schneiden.
            
            \item Sojasauce und Erdnussbutter in die Pfanne geben, etwas Wasser hinzufügen und zu einer halbflüssigen Sauce verrühren. Reisweinessig und Sesamöl unterrühren.
            
            \item Nudeln zur Sauce geben, Frühlingszwiebeln und Sesam hinzufügen und alles gut vermengen.
        \end{enumerate}
    }
    
    \suggestion{%
        Wenn du es schärfer magst, einfach mehr Chiliöl rein. Mit Limettensaft wird es frischer.
    }
    
    \hint{%
        Falls die Sauce zu dick wird, einfach schluckweise Nudelwasser hinzufügen – das macht sie schön cremig.
    }
    
\end{recipe}